\usepackage{amsmath,amssymb,mathrsfs,fancyhdr,enumerate,multirow,makecell,currfile,venndiagram,fontawesome,pifont,pgf-pie,pgfplots,yhmath,twemojis}
\usepackage{zref-lastpage}
\usepackage{zref-user}
\usepackage{mathpazo}
\usepackage{tikz-3dplot,tikz,tkz-tab,tabvar,tkz-euclide}
\usetikzlibrary{arrows,calc,intersections,angles,snakes,quotes,backgrounds,shapes.geometric,patterns, shapes.symbols}
\usepackage[loigiai]{ex_test}
%\usepackage[dethi]{ex_test}
\pgfplotsset{compat=1.9}
\usepgfplotslibrary{fillbetween}
\usepackage[tikz]{bclogo}
\usepackage[
top=1.3cm,
bottom=1.3cm,
left=1.5cm,
right=1cm,
headsep=10pt,   % nhỏ lại
footskip=22pt  % khoảng cách từ nội dung xuống footer
]{geometry}
\usepackage[hidelinks,unicode,hypertexnames=false]{hyperref}
\usepackage[final]{pdfpages}
\usepackage{tcolorbox,diagbox}
\tcbuselibrary{theorems, breakable, skins}
\usepackage{xcolor}
%=====
\setcounter{secnumdepth}{4}
\setcounter{tocdepth}{1}
\usepackage{titlesec}
\usepackage{titletoc}
\makeatletter
%%%%%%%%%%%%%%======Khai báo màu
\definecolor{myred}{RGB}{210,0,0}
%Khai báo chapter, section, subsection, subsubsection
% chapter
\titleformat{\chapter}
{\bfseries\large\color{myred}}
{CHƯƠNG \thechapter.}
{0.5em}{\MakeUppercase}
[\tikz{\draw[myred, line width=1pt,double,thick] (0,0)--(\textwidth,0);}]
\titlespacing*{\chapter}{0pt}{-5pt}{5pt}
% Section
\titleformat{\section}
{\bfseries\large\color{myred}}
{BÀI \arabic{section}.}{0.3em}{\MakeUppercase}
\titlespacing*{\section}{0pt}{3pt}{5pt}
% Subsection
\renewcommand{\thesubsection}{\Alph{subsection}.}
\titleformat{\subsection}
{\bfseries}
{\Alph{subsection}.}{0.3em}{}
\titlespacing*{\subsection}{0pt}{3pt}{5pt}
% Subsubsection
\titleformat{\subsubsection}
{\bfseries}
{\arabic{subsubsection}.}{0.25em}{}
\titlespacing*{\subsubsection}{0pt}{3pt}{0pt}
% Paragraph
\renewcommand{\theparagraph}{\arabic{paragraph})}
\titleformat{\paragraph}
{\bfseries}
{\arabic{subsubsection}.\arabic{paragraph}.}{0.3em}{}
\titlespacing*{\paragraph}{0pt}{3pt}{0pt}
%part
\titleformat{\part}[display]
{\Large\bfseries\color{red!50!blue}}{
	\thispagestyle{empty}
	\centering
	\begin{tikzpicture}
		\fill[red!50!blue] (-1.5,-1) rectangle (1.5,1);
		\draw(0,1.5)node{\fontfamily{ccr}\selectfont\Large PHẦN};
		\draw(0,0)node{\color{white}\fontsize{60pt}{1pt}\selectfont\thepart};
	\end{tikzpicture}
}
{0.5em}
{\titlerule[2.5pt]\vskip3pt\titlerule\vskip4pt\centering\fontsize{60pt}{1pt}\fontfamily{qhv}\selectfont\MakeUppercase}
\renewcommand\subsection{\@startsection{subsection}{2}{\z@}%
	{1ex\@plus -1ex \@minus -.2ex}%
	{0pt}%
	{\normalfont\normalsize\bfseries}}
\makeatother
%=====Đổi màu phương án đúng
\def\colorEX{\color{blue}}
%Môi trường bài tập tự luận
\newtheorem{bt}{Bài}
%Định nghĩa thêm cách gõ vectơ, khuyến khích gõ \overrightarrow
\usepackage{esvect}
\def\vec{\vv}
\def\overrightarrow{\vv}
%Làm lại mục lục part
\makeatletter
\renewcommand*\l@part[2]{%
	\ifnum \c@tocdepth >\m@ne
	\addpenalty{-\@highpenalty}%
	\vskip 1.0em \@plus\p@
	\setlength\@tempdima{1.5em}%
	\begingroup
	\parindent \z@ \rightskip \@pnumwidth
	\parfillskip -\@pnumwidth
	\leavevmode
	\advance\leftskip\@tempdima
	\hskip -\leftskip
	\colorbox{violet}{\strut%
		\makebox[\dimexpr\textwidth-1\fboxsep-6pt\relax][l]{%
			%\fontsize{12pt}{1pt}
			\color{white}\bfseries\fontfamily{qag}\selectfont\hspace*{0.3\textwidth} PHẦN #1%
			\nobreak\hfill\nobreak\hb@xt@\@pnumwidth{\hss #2}}}\par\smallskip
	\penalty\@highpenalty
	\endgroup
	\fi}
\makeatother
%Lệnh song song
\DeclareSymbolFont{symbolsC}{U}{txsyc}{m}{n}
\DeclareMathSymbol{\varparallel}{\mathrel}{symbolsC}{9}
\DeclareMathSymbol{\parallel}{\mathrel}{symbolsC}{9}
%Hệ hoặc, hệ và
\newcommand{\hoac}[1]{
	\left[\begin{aligned}#1\end{aligned}\right.}
\newcommand{\heva}[1]{
	\left\{\begin{aligned}#1\end{aligned}\right.}
%Giãn dòng
% \renewcommand{\baselinestretch}{1.4}
%Thay đổi đếm alpha môi trường enumerate
\makeatletter
\renewcommand{\labelenumi}{\alph{enumi})}
\makeatother
%Tạo biến đếm số đề và khung tiêu đề
\def\khoi{12}
\def\made{1234}
\newcounter{demsode}[chapter]
\setcounter{demsode}{1}
\newenvironment{name}[8]{
	\setcounter{phan}{0}
	\setcounter{page}{1}
	\noindent
	\begin{minipage}[]{0.45\textwidth}
		\centering
		{#1}\\
		{#2}\\
		{#3}\\
		{#4}
	\end{minipage}
	\begin{minipage}[]{0.55\textwidth}
		\centering
		{#5}\\
		{#6}\\
		{#7}\\
		{#8}
	\end{minipage}
	\\[1mm]
	\tikz \draw[double] (0,0)--(\textwidth,0);\\
	\textit{Họ và tên thí sinh:} \tikz \draw[dotted,very thin] (0,0)--(\textwidth/3,0); 
	\textit{Số báo danh:}\tikz \draw[dotted,very thin] (0,0)--(\textwidth/3.8,0);\quad
	\fbox{\color{blue} \bf \large\fontfamily{anttlc}\selectfont \made}
	%=====
	\addcontentsline{toc}{section}{Đề \thedemsode. Đề KSCL năm học 2025-2026, #2, Môn Toán \khoi}
	\addtocounter{demsode}{1}
	\gdef\socaulc{12}
	\gdef\socauds{4}
	\gdef\socaukq{6}
	\gdef\socautl{3}
}

%%%%%%%%%%%%%%
\newcommand{\indapan}[1]{
%	\begin{center}
%		\fbox{\large \textbf{\color{purple} \large BẢNG ĐÁP ÁN}}
%	\end{center}
	\IfFileExists{#1-Phan-I}{
		\noindent
		\inputansbox{12}{#1-Phan-I}
	}{\null}
	\IfFileExists{#1-Phan-II}{
		\noindent
		\inputansbox{2}{#1-Phan-II}
	}{\null}
	\IfFileExists{#1-Phan-III}{
		\noindent
		\inputansbox{3}{#1-Phan-III}
	}{\null}
}
\newcommand{\indapancuoi}[1]{
	\phantomsection
	\begin{center}
		\fbox{\large \textbf{\color{purple} \large BẢNG ĐÁP ÁN ĐỀ SỐ \thedemsode}}
	\end{center}
	\IfFileExists{#1-Phan-I}{
		\noindent
		\inputansbox{12}{#1-Phan-I}
	}{\null}
	\IfFileExists{#1-Phan-II}{
		\noindent
		\inputansbox[2]{2}{#1-Phan-II}
	}{\null}
	\IfFileExists{#1-Phan-III}{
		\noindent
		\inputansbox[2]{6}{#1-Phan-III}
	}{\null}
	\addcontentsline{toc}{section}{Đề số \thedemsode}
	\addtocounter{demsode}{1}
	\newpage}
%Cấu trúc đề mới
%\def\socaulc{12}
%\def\socauds{2}
%\def\socaukq{4}
%\def\socautl{3}
\newcounter{phan}
\setcounter{phan}{0}

\renewcommand{\thephan}{\Roman{phan}}

% PHẦN I
\newcommand{\caulc}{
	\stepcounter{phan}
	\noindent\textbf{PHẦN \thephan. Câu trắc nghiệm nhiều phương án lựa chọn.}
	\setcounter{ex}{0}
}

% PHẦN II
\newcommand{\cauds}{
	\stepcounter{phan}
	\noindent\textbf{PHẦN \thephan. Câu trắc nghiệm đúng sai.}
	\setcounter{ex}{0}
}

% PHẦN III
\newcommand{\caukq}{
	\stepcounter{phan}
	\noindent\textbf{PHẦN \thephan. Câu trắc nghiệm trả lời ngắn.}
	\setcounter{ex}{0}
}

% PHẦN IV
\newcommand{\cautl}{
	\stepcounter{phan}
	\noindent\textbf{PHẦN \thephan. Câu tự luận.}
	\setcounter{ex}{0}
	\setcounter{bt}{0}
}
%Định nghĩa Head và Foot
\pagestyle{fancy}
\fancyhf{}
% Header
%\fancyhead[L]{\color{red!50!white}\faPhone\,0812.8787.92}
%\fancyhead[R]{\color{red!50!white}Lưu hành nội bộ}
% Footer
%\fancyfoot[L]{\color{red!50!white}Ths: Hồ Ngọc Nhất Linh}
%\fancyfoot[R]{\color{red!50!white}Ths: Hồ Ngọc Nhất Linh}
% Số trang:
% LE = trang chẵn bên trái
% RO = trang lẻ bên phải
\fancyfoot[LE,RO]{Trang \thepage/\zpageref{\made-lastpage} \,--\, Mã đề \made}
% Quan trọng: sửa style plain
\fancypagestyle{plain}{
%	\fancyhf{}
%	\fancyhead[L]{\color{red!50!white}\faPhone\,0812.8787.92}
%	\fancyhead[R]{\color{red!50!white}Lưu hành nội bộ}
%	\fancyfoot[L]{\color{red!50!white}Ths: Hồ Ngọc Nhất Linh}
%	\fancyfoot[R]{\color{red!50!white}Ths: Hồ Ngọc Nhất Linh}
	\fancyfoot[LE,RO]{\thepage/\made}
}
\renewcommand{\headrulewidth}{0pt}
%\renewcommand{\footrulewidth}{0.5pt}
%\makeatletter
%\renewcommand{\headrule}{%
%	{\color{red}\hrule \@height \headrulewidth \@width \headwidth}
%}
%\renewcommand{\footrule}{%
%	{\color{red}\hrule \@height \footrulewidth \@width \headwidth}
%}
%\makeatother
%=====
%%%%%%%%%===Tạo khung
\newtcolorbox{khung}{
	before upper=\raggedright,
	colback=white,    % màu nền
	boxrule=0.5pt,
	colframe=gray!50!black,  % màu khung
%	title=\textbf{Định nghĩa},
	fonttitle=\bfseries,
	breakable,               % cho phép xuống dòng qua trang
	arc=1mm,                  % bo góc
	left=2pt,
	right=2pt,
	top=2pt,
	bottom=2pt
}
%%%%%%%%%%---Khai báo định nghĩa, định lý, lưu ý, chú ý, note
\newenvironment{dn}
{\begin{khung}\textcolor{red}{\textbf{Định nghĩa. }}}{\end{khung}}
%%%%%%%%%%%%%%%%%%%%%%%%%%%%%%%%%%%%%%%%%%%%%%%%%%%%%%%%%%%%%%%%%%
\newenvironment{dl}
{\begin{khung}\textcolor{blue}{\textbf{Định lý. }}}{\end{khung}}
%%%%%%%%%%%%%%%%%%%%%%%%%%%%%%%%%%%%%%%%%%%%%%%%%%%%%%%%%%%%%%%%%%
\newenvironment{luuy}
{\begin{khung}\textcolor{purple}{\textbf{Lưu ý. }}}{\end{khung}}
%%%%%%%%%%%%%%%%%%%%%%%%%%%%%%%%%%%%%%%%%%%%%%%%%%%%%%%%%%%%%%%%%%
\newenvironment{nx}
{\begin{khung}\textcolor{purple}{\textbf{Nhận xét. }}}{\end{khung}}
%%%%%%%%%%%%%%%%%%%%%%%%%%%%%%%%%%%%%%%%%%%%%%%%%%%%%%%%%%%%%%%%%%
\newenvironment{pp}
{\begin{khung}\textcolor{red}{\textbf{Phương pháp. }}}{\end{khung}}
%%%%%%%%%%%%%%%%%%%%%%%%%%%%%%%%%%%%%%%%%%%%%%%%%%%%%%%%%%%%%%%%%%
\newenvironment{chuy}
{\begin{khung}\textcolor{green}{\textbf{Chú ý. }}}{\end{khung}}
%%%%%%%%%%%%%%%%%%%%%%%%%%%%%%%%%%%%%%%%%%%%%%%%%%%%%%%%%%%%%%%%%%
\newenvironment{note}
{\begin{khung}\textcolor{cyan}{\textbf{Ghi chú. }}}{\end{khung}}
%%%%%%%%%%%%%%%%%%%%%%%%%%%%%%%%%%%%%%%%%%%%%%%%%%%%%%%%%%%%%%%%%%
\newenvironment{tc}
{\begin{khung}\textcolor{green!70!black}{\textbf{Tính chất. }}}{\end{khung}}
%%%%%%%%%%%%%%%%%%%%%%%%%%%%%%%%%%%%%%%%%%%%%%%%%%%%%%%%%%%%%%%%%%
\newcounter{vd}
\newenvironment{vd}
{\refstepcounter{vd}%
\begin{khung}\textcolor{black}{\textbf{Ví dụ \thevd. }}}{\end{khung}}
%%%%%%%%%%%%%%%%%%%%%%%%%%%%%%%%%%%%%%%%%%%%%%%%%%%%%%%%%%%%%%%%%%
%%%%%%%%%%%%%%%======dòng chấm, oli
\newcommand{\dongcham}[1]{
	\def\sod{#1}
	\pgfmathsetmacro{\sodong}{2*\sod -1} 
	\columnsep=10pt
	\vspace*{-3.5mm}
	\begin{multicols}{2}
		\foreach \dotline in{1,...,\sodong}
		{\noindent\color{black!90}{\tiny\dotfill}\\[2mm]}
		{\noindent\color{black!90}{\tiny\dotfill}}\\[-4mm]
	\end{multicols}
}
\newcommand{\dc}[2][1.3]{%
	\begingroup
	% đặt parskip = 0 chỉ trong nhóm này
	\setlength{\parskip}{0pt}%
	\setlength{\parindent}{0pt}% nếu bạn không muốn indent
	% bắt đầu in dòng dotfill
	\par\nobreak
	\noindent\rule{0pt}{#1\baselineskip}%
	\foreach \i in {1,...,#2}{%
		\noindent
		\makebox[\linewidth]{\rule{0pt}{#1\baselineskip}\tiny\dotfill}\par
	}%
	\endgroup
}
%%%%%%%%%%%%%%%%%%%%%%%%%%%%%%%%%%%%%%
\newcommand{\oli}[1]{%
	\noindent
	\ifdim\textwidth>10cm
	\tikz{\draw[gray!70, dashed, step=\textwidth/24] 
		(0,0) grid (\textwidth,{#1*\textwidth/24});}
	\else
	\tikz{\draw[gray!70, dashed, step=\textwidth/13] 
		(0,0) grid (\textwidth,{#1*\textwidth/13});}
	\fi
}
%%%%%%==============lệnh chia hai cột
\newcommand{\chiahaicot}[2]{
	\noindent
	\begin{minipage}[t]{0.5\textwidth}
		\begin{khung}
			#1
		\end{khung}
	\end{minipage}
	\begin{minipage}[t]{0.5\textwidth}
		\begin{khung}
			#2
		\end{khung}
	\end{minipage}
}
%%%%%%%%%%%%======================
\newcommand{\anloigiai}{
	\hideansEX{vd}
	\hideansEX{ex}
	\hideansEX{bt}
}
\newcommand{\tatdongcham}{
	\renewcommand{\dc}[2][1.2]{}
	\renewcommand{\dongcham}[1]{}
}